% !LW recipe = pdflatex

\documentclass[12pt,a4paper]{article}

% Permite escrever caracteres acentuados.
\usepackage[utf8]{inputenc}

% Melhora a codificação das fontes no PDF.
\usepackage[T1]{fontenc}

% Regras de idioma do português europeu.
\usepackage[portuguese]{babel}

% Pacotes matemáticos.
\usepackage{amsmath}
\usepackage{amssymb}
\usepackage{amsthm}

% Configuração das margens.
\usepackage{geometry}
\geometry{
    top=2.5cm,
    bottom=2.5cm,
    left=2.5cm,
    right=2.5cm
}

% Símbolos das operações que serão construídas.
\newcommand{\mov}{\mathbin{\#}}
\newcommand{\rep}{\mathbin{@}}

\title{Matemática como Movimento}
\author{Timóteo Cruz}
\date{\today}

\begin{document}

\maketitle

\section{Introdução}

\subsection{A ideia fundamental}

Começamos com o conjunto dos números naturais:

\[
    \mathbb{N}=\{1,2,3,4,5,\ldots\}.
\]

Os números naturais surgiram da necessidade de contar objetos do
quotidiano. Para determinar quantos animais possuímos, por exemplo,
podemos associar cada animal a um dedo, a uma pedra ou a uma marca.

Essa é uma correspondência um a um: cada objeto contado recebe
exatamente uma marca.

Ao escrevermos os números naturais na ordem

\[
    1,2,3,4,5,\ldots,
\]

já reconhecemos uma primeira estrutura: depois de cada número aparece
um número seguinte.

Neste primeiro momento, porém, não assumiremos que conhecemos operações
como a soma ou a multiplicação. Queremos observar como essas operações
podem surgir a partir de perguntas simples.

Como os números estão organizados numa sequência, podemos imaginar que
nos movemos sobre eles.

A nossa primeira pergunta será:

\begin{center}
    \emph{Como podemos descrever movimentos sobre os números naturais?}
\end{center}

Consideremos o trecho

\[
    A=\{1,2,3,4,5,6\}\subseteq\mathbb{N}.
\]

Imaginemos agora que partimos de uma origem ainda não nomeada e damos
passos de comprimento \(2\).

Depois do primeiro passo, alcançamos o número \(2\). Depois do segundo,
alcançamos o número \(4\). Depois do terceiro, alcançamos o número \(6\).

Podemos representar a composição desses três movimentos por

\[
    2\mov 2\mov 2=6.
\]

Aqui, cada ocorrência do número \(2\) representa um passo de comprimento
\(2\). O primeiro \(2\) não representa a posição inicial: ele representa
o primeiro movimento.

O símbolo \(\mov\) será, portanto, a nossa primeira operação. Neste
momento, ele representa a composição de dois movimentos.

Mais tarde investigaremos quais propriedades essa operação possui.

\subsection{Repetição de um mesmo passo}

No movimento anterior aparece uma estrutura interessante:

\begin{itemize}
    \item o comprimento de cada passo é \(2\);
    \item o mesmo passo aparece \(3\) vezes;
    \item o ponto alcançado é \(6\).
\end{itemize}

A repetição do mesmo movimento sugere a criação de uma nova notação.

Definimos

\[
    a\rep n
    =
    \underbrace{a\mov a\mov\cdots\mov a}_{n\text{ ocorrências de }a}.
\]

Assim, \(a\rep n\) significa repetir \(n\) vezes um movimento de
comprimento \(a\).

No nosso primeiro exemplo,

\[
    2\rep 3=2\mov 2\mov 2=6.
\]

Como também podemos representar o número \(3\) por

\[
    1\mov 1\mov 1=3,
\]

obtemos

\[
    2\rep(1\mov 1\mov 1)
    =
    2\rep 3
    =
    6.
\]

Outro exemplo seria

\[
    4\rep 3
    =
    4\mov 4\mov 4
    =
    12.
\]

Neste estágio, \(\rep\) é apenas uma operação de repetição. Ainda
precisamos investigar as suas propriedades antes de identificá-la
com aquilo que normalmente chamamos de multiplicação.

\subsection{Duas maneiras pelas quais uma operação pode surgir}

Podemos agora distinguir duas maneiras pelas quais uma operação pode
aparecer.

A operação \(\mov\) surgiu da necessidade de representar a composição
de movimentos.

A operação \(\rep\), por outro lado, surgiu quando reconhecemos uma
repetição em expressões construídas com \(\mov\).

Por exemplo, considere a expressão

\[
    2\mov 2\mov 2\mov 2\mov 2.
\]

Em vez de escrever cinco vezes o mesmo passo, podemos comprimir a
expressão:

\[
    2\rep 5.
\]

Portanto, a operação \(\rep\) não aparece inicialmente como uma
operação completamente independente. Ela surge como uma maneira de
reconhecer, nomear e comprimir uma regularidade presente na repetição
da operação \(\mov\).

Somente depois de estudarmos as propriedades dessa nova operação
poderemos reconhecer que ela corresponde à multiplicação usual dos
números naturais.
\subsection{Primeiras propriedades da operação \texorpdfstring{$\mov$}{#}}

Até agora, observámos alguns movimentos particulares. Contudo, verificar
uma propriedade em alguns exemplos não significa que ela seja verdadeira
para todos os números naturais.

Por exemplo, podemos calcular

\[
2\mov3=5
\]

e

\[
3\mov2=5,
\]

mas isso, por si só, não demonstra que

\[
a\mov b=b\mov a
\]

para quaisquer números naturais (a) e (b).

Surge, portanto, uma nova necessidade:

\begin{center}
\emph{Como podemos provar uma afirmação para todos os números naturais?}
\end{center}

\subsubsection{O número seguinte e a definição precisa de movimento}

Representaremos por (S(n)) o número natural que aparece imediatamente
depois de (n). Assim,

\[
S(0)=1,\qquad S(1)=2,\qquad S(2)=3,
\]

e assim sucessivamente.

Podemos agora tornar a operação $\mov$ mais precisa.

Definimos:

\[
a\mov0=a
\]

e

\[
a\mov S(b)=S(a\mov b).
\]

A primeira igualdade diz que realizar zero movimentos a partir de (a)
nos deixa em (a).

A segunda diz que realizar (S(b)) movimentos significa realizar (b)
movimentos e depois mais um movimento unitário.

Por exemplo,

\[
\begin{aligned}
2\mov3
&=2\mov S(2)\
&=S(2\mov2)\
&=S(S(2\mov1))\
&=S(S(S(2\mov0)))\
&=S(S(S(2)))\
&=5.
\end{aligned}
\]

Dessa maneira, $a\mov b$ representa o ponto alcançado quando partimos
de (a) e avançamos (b) passos unitários.

\subsubsection{Uma nova ferramenta: a indução}

Para provar que uma propriedade vale para todos os números naturais,
podemos seguir dois passos:

\begin{enumerate}
\item provar que a propriedade vale para (0);
\item provar que, sempre que ela vale para um número (n), também
vale para o número seguinte (S(n)).
\end{enumerate}

Esse procedimento é chamado de \emph{indução matemática}.

A ideia é semelhante a um movimento sobre os próprios números naturais.
Começamos no (0) e mostramos que a verdade pode passar de cada número
para o número seguinte.

\subsubsection{A ordem dos movimentos importa?}

Comparemos:

\[
2\mov3=5
\]

e

\[
3\mov2=5.
\]

Ambos os movimentos chegam ao mesmo número. Isso sugere que

\[
a\mov b=b\mov a.
\]

Essa propriedade é chamada de \emph{comutatividade}.

Antes de demonstrá-la, precisamos de dois resultados auxiliares.

\paragraph{Primeiro resultado.}

Para qualquer número natural (a),

\[
0\mov a=a.
\]

\begin{proof}
Faremos indução sobre (a).

Para (a=0),

\[
0\mov0=0.
\]

Suponhamos agora que

\[
0\mov a=a.
\]

Então,

\[
\begin{aligned}
0\mov S(a)
&=S(0\mov a)\
&=S(a).
\end{aligned}
\]

Logo, a propriedade passa de (a) para (S(a)). Portanto,

\[
0\mov a=a
\]

para todo número natural (a).
\end{proof}

\paragraph{Segundo resultado.}

Para quaisquer números naturais (a) e (b),

\[
S(a)\mov b=S(a\mov b).
\]

\begin{proof}
Faremos indução sobre (b).

Para (b=0),

\[
S(a)\mov0=S(a)=S(a\mov0).
\]

Suponhamos que

\[
S(a)\mov b=S(a\mov b).
\]

Então,

\[
\begin{aligned}
S(a)\mov S(b)
&=S\bigl(S(a)\mov b\bigr)\
&=S\bigl(S(a\mov b)\bigr)\
&=S\bigl(a\mov S(b)\bigr).
\end{aligned}
\]

Portanto, o resultado vale para todo (b).
\end{proof}

Podemos agora demonstrar a comutatividade.

\begin{proof}
Fixemos um número natural (a) e façamos indução sobre (b).

Para (b=0),

\[
a\mov0=a=0\mov a.
\]

Suponhamos que

\[
a\mov b=b\mov a.
\]

Então,

\[
\begin{aligned}
a\mov S(b)
&=S(a\mov b)\
&=S(b\mov a)\
&=S(b)\mov a.
\end{aligned}
\]

Logo,

\[
a\mov b=b\mov a
\]

para todos os números naturais (a) e (b).
\end{proof}

A ordem dos movimentos, portanto, não altera o ponto final.

Isso não significa que todas as operações sejam comutativas. Podemos,
por exemplo, inventar uma operação $\triangleright$ definida por

\[
a\triangleright b=a.
\]

Essa operação simplesmente conserva o primeiro número. Assim,

\[
2\triangleright3=2,
\]

enquanto

\[
3\triangleright2=3.
\]

Logo, a ordem dos elementos pode alterar o resultado de uma operação.

\subsubsection{A maneira de agrupar importa?}

Comparemos:

\[
\begin{aligned}
(2\mov3)\mov4
&=5\mov4\
&=9
\end{aligned}
\]

e

\[
\begin{aligned}
2\mov(3\mov4)
&=2\mov7\
&=9.
\end{aligned}
\]

Isso sugere que

\[
(a\mov b)\mov c=a\mov(b\mov c).
\]

Essa propriedade é chamada de \emph{associatividade}.

Ela afirma que, quando realizamos vários movimentos sucessivos, a
maneira de colocar os parênteses não altera o destino.

\begin{proof}
Fixemos (a) e (b) e façamos indução sobre (c).

Para (c=0),

\[
\begin{aligned}
(a\mov b)\mov0
&=a\mov b\
&=a\mov(b\mov0).
\end{aligned}
\]

Suponhamos que

\[
(a\mov b)\mov c=a\mov(b\mov c).
\]

Então,

\[
\begin{aligned}
(a\mov b)\mov S(c)
&=S\bigl((a\mov b)\mov c\bigr)\
&=S\bigl(a\mov(b\mov c)\bigr)\
&=a\mov S(b\mov c)\
&=a\mov\bigl(b\mov S(c)\bigr).
\end{aligned}
\]

Portanto,

\[
(a\mov b)\mov c=a\mov(b\mov c)
\]

para todos os números naturais (a), (b) e (c).
\end{proof}

A associatividade e a comutatividade são propriedades diferentes. Uma
operação pode possuir uma delas sem possuir a outra.

\paragraph{Uma operação associativa, mas não comutativa.}

Consideremos novamente

\[
a\triangleright b=a.
\]

Temos

\[
\begin{aligned}
(a\triangleright b)\triangleright c
&=a\triangleright c\
&=a
\end{aligned}
\]

e

\[
\begin{aligned}
a\triangleright(b\triangleright c)
&=a\triangleright b\
&=a.
\end{aligned}
\]

Logo, $\triangleright$ é associativa.

Entretanto,

\[
a\triangleright b=a
\]

e

\[
b\triangleright a=b,
\]

que geralmente são diferentes. Portanto, ela não é comutativa.

\paragraph{Uma operação comutativa, mas não associativa.}

Definamos $a\diamond b$ como a distância entre as posições (a) e
(b) na sequência dos números naturais.

A distância entre (a) e (b) é igual à distância entre (b) e (a).
Portanto,

\[
a\diamond b=b\diamond a.
\]

Entretanto,

\[
\begin{aligned}
(1\diamond2)\diamond3
&=1\diamond3\
&=2,
\end{aligned}
\]

enquanto

\[
\begin{aligned}
1\diamond(2\diamond3)
&=1\diamond1\
&=0.
\end{aligned}
\]

Assim,

\[
(1\diamond2)\diamond3
\neq
1\diamond(2\diamond3).
\]

Logo, $\diamond$ é comutativa, mas não associativa.

\subsubsection{Existe algum movimento que não altera nada?}

Procuramos um elemento (e) tal que

\[
a\mov e=a
\]

para todo número natural (a).

Pela própria definição da operação,

\[
a\mov0=a.
\]

Além disso, pela comutatividade,

\[
0\mov a=a.
\]

Portanto, o zero é o $\emph{elemento neutro}$ da operação $\mov$.

Inicialmente, o zero apareceu como a ausência de movimento. Entretanto,
ele não permanece apenas como uma ideia de ausência. Ele pode participar
de operações, ser comparado com outros números e ocupar uma posição
definida na estrutura que estamos a construir.

Nesse sentido, a ausência de movimento passa a ser representada por um
novo número.

Como começámos com

\[
\mathbb{N}={1,2,3,\ldots},
\]

a descoberta do zero obriga-nos a ampliar o nosso conjunto:

\[
\mathbb{N}_0={0,1,2,3,\ldots}.
\]

A partir deste momento, trabalharemos em $\mathbb{N}_0$.

\subsection{Primeiras propriedades da operação \texorpdfstring{$\rep$}{@}}

Recordemos que

\[
a\rep n
=
\underbrace{a\mov a\mov\cdots\mov a}
_{n\text{ ocorrências de }a}.
\]

Podemos também definir essa operação recursivamente:

\[
a\rep0=0
\]

e

\[
a\rep S(n)=(a\rep n)\mov a.
\]

A primeira igualdade diz que repetir um movimento zero vezes não produz
movimento algum.

A segunda diz que repetir (a) por (S(n)) vezes significa repetir
(a) por (n) vezes e depois acrescentar mais uma ocorrência de (a).

\subsubsection{Qual é o elemento neutro da repetição?}

Como repetir (a) uma única vez produz o próprio (a),

\[
a\rep1=a.
\]

Também temos

\[
1\rep a
=
\underbrace{1\mov1\mov\cdots\mov1}
_{a\text{ ocorrências de }1}
=
a.
\]

Portanto, o elemento neutro de $\rep$ é o número (1).

O zero possui outro comportamento:

\[
a\rep0=0
\]

e

\[
0\rep a=0.
\]

Assim, o zero não é neutro para a repetição. Ele é um elemento
\emph{absorvente}: quando participa de uma repetição, o resultado é
zero.

\subsubsection{A repetição é comutativa?}

Comparemos:

\[
2\rep3=2\mov2\mov2=6
\]

e

\[
3\rep2=3\mov3=6.
\]

Isso sugere que

\[
a\rep b=b\rep a.
\]

Podemos imaginar $a\rep b$ como uma organização retangular formada
por (b) grupos contendo (a) objetos em cada grupo.

Se observarmos o mesmo retângulo na outra direção, veremos (a) grupos
contendo (b) objetos em cada grupo.

O número total de objetos não se altera. Essa observação conduz à
comutatividade da repetição:

\[
a\rep b=b\rep a.
\]

Uma demonstração completamente formal também pode ser feita por
indução.

\subsubsection{A repetição é associativa?}

Consideremos

\[
(a\rep b)\rep c.
\]

Primeiro, repetimos (a) por (b) vezes. Depois, repetimos todo esse
bloco por (c) vezes.

No total, temos (c) blocos, cada um contendo (b) ocorrências de
(a). Isso equivale a repetir (a) por $b\rep c$ vezes.

Portanto,

\[
(a\rep b)\rep c=a\rep(b\rep c).
\]

A operação $\rep$ também é associativa.

\subsection{Distribuição de movimentos}

Consideremos

\[
2\rep(3\mov4).
\]

Como

\[
3\mov4=7,
\]

essa expressão significa repetir o movimento de comprimento (2) por
sete vezes:

\[
2\rep7
=
2\mov2\mov2\mov2\mov2\mov2\mov2.
\]

Podemos separar essas sete ocorrências em dois grupos: um grupo com três
ocorrências e outro com quatro ocorrências:

\[
\begin{aligned}
2\rep(3\mov4)
&=
(2\mov2\mov2)
\mov
(2\mov2\mov2\mov2)\
&=
(2\rep3)\mov(2\rep4).
\end{aligned}
\]

Isso sugere a propriedade

\[
a\rep(b\mov c)
=
(a\rep b)\mov(a\rep c).
\]

Essa propriedade é chamada de \emph{distributividade de $\rep$ sobre
$\mov$}.

A razão pela qual podemos distribuir a repetição é que repetir (a)
por $b\mov c$ vezes produz uma sequência de repetições que pode ser
separada em duas partes:

\[
\underbrace{a\mov\cdots\mov a}
_{b\text{ ocorrências}}
\mov
\underbrace{a\mov\cdots\mov a}
_{c\text{ ocorrências}}.
\]

Também podemos começar com uma expressão como

\[
(3\mov4)\rep2.
\]

Agora estamos a repetir o movimento $3\mov4$ duas vezes:

\[
(3\mov4)\mov(3\mov4).
\]

Usando a associatividade e a comutatividade de $\mov$, podemos
reorganizar os movimentos:

\[
\begin{aligned}
(3\mov4)\mov(3\mov4)
&=3\mov4\mov3\mov4\
&=3\mov3\mov4\mov4\
&=(3\rep2)\mov(4\rep2).
\end{aligned}
\]

Assim, também temos

\[
(a\mov b)\rep c
=
(a\rep c)\mov(b\rep c).
\]

\subsubsection{A operação \texorpdfstring{$\mov$}{#} é distributiva
sobre \texorpdfstring{$\rep$}{@}?}

Para que $\mov$ fosse distributiva sobre $\rep$, precisaríamos ter

\[
a\mov(b\rep c)
=
(a\mov b)\rep(a\mov c)
\]

para quaisquer (a), (b) e (c).

Tomemos (a=2), (b=3) e (c=4).

O primeiro lado é

\[
\begin{aligned}
2\mov(3\rep4)
&=2\mov12\
&=14.
\end{aligned}
\]

O segundo lado é

\[
\begin{aligned}
(2\mov3)\rep(2\mov4)
&=5\rep6\
&=30.
\end{aligned}
\]

Como

\[
14\neq30,
\]

a operação $\mov$ não é distributiva sobre $\rep$.

Temos, assim, uma situação assimétrica:

\[
\rep\text{ é distributiva sobre }\mov,
\]

mas

\[
\mov\text{ não é distributiva sobre }\rep.
\]

\subsection{Reconhecimento das operações}

A operação $\mov$ foi definida como o avanço de um determinado número
de passos unitários. Ela produz exatamente os mesmos resultados que a
operação normalmente chamada de soma.

A operação $\rep$ foi definida como a repetição de um mesmo movimento.
Ela produz exatamente os mesmos resultados que a operação normalmente
chamada de multiplicação.

Portanto, podemos identificar

\[
a\mov b=a+b
\]

e

\[
a\rep b=a\cdot b.
\]

Não fazemos essa identificação apenas porque as operações possuem
propriedades semelhantes. Fazemo-la porque as próprias definições
produzem, para todos os números naturais, os mesmos resultados da soma
e da multiplicação usuais.

A partir deste ponto, usaremos as notações conhecidas (+) e
$\cdot$. Mais adiante, voltaremos a procurar novas operações e novas
formas de movimento sobre os números.

\section{O surgimento das incógnitas}

\subsection{Um movimento desconhecido}

Consideremos a igualdade

\[
2+x=6.
\]

Aqui sabemos:

\begin{itemize}
\item o ponto de onde partimos;
\item o ponto ao qual queremos chegar;
\item mas não sabemos qual movimento deve ser realizado.
\end{itemize}

O símbolo (x) representa esse movimento desconhecido.

Uma letra usada para representar uma quantidade cujo valor desejamos
determinar é chamada de \emph{incógnita}.

A pergunta que precisamos responder é:

\begin{center}
\emph{Qual movimento deve ser realizado para alcançar um destino
previamente escolhido?}
\end{center}

Essa pergunta conduz ao conceito de equação.

Uma \emph{equação} é uma igualdade que contém uma ou mais incógnitas. Uma
solução da equação é um valor que, ao ser colocado no lugar da incógnita,
torna a igualdade verdadeira.

No nosso exemplo,

\[
2+x=6,
\]

o valor (x=4) é uma solução, pois

\[
2+4=6.
\]

O conjunto no qual procuramos as soluções também é importante. Uma mesma
equação pode ter solução num conjunto numérico e não ter solução noutro.

\subsection{Como resolver uma equação?}

Na equação

\[
2+x=6,
\]

procuramos o movimento que, partindo de (2), nos leva até (6).

Por inspeção dos números naturais, percebemos que esse movimento é (4).
Entretanto, nem sempre será fácil encontrar uma solução experimentando
número por número.

Isso sugere a necessidade de uma nova operação.

Queremos uma operação que responda à seguinte pergunta:

\begin{center}
\emph{Que movimento transforma o ponto inicial no ponto final?}
\end{center}

Introduziremos temporariamente o símbolo $(\ominus)$.

Definimos

\[
a\ominus b=x
\]

quando

\[
b+x=a.
\]

Portanto, $a\ominus b$ representa o movimento que, partindo de (b),
nos leva até (a).

No exemplo anterior,

\[
x=6\ominus2,
\]

pois procuramos o movimento que, somado a (2), produz (6):

\[
2+(6\ominus2)=6.
\]

Assim,

\[
6\ominus2=4.
\]

Essa nova operação não é apenas uma forma diferente de soma. Ela resolve
o problema inverso da soma.

Na soma, conhecemos os movimentos e procuramos o destino:

\[
2+4=6.
\]

Na nova operação, conhecemos o ponto inicial e o destino, mas procuramos
o movimento:

\[
6\ominus2=4.
\]

\subsection{Subtração e distância}

É importante distinguir a nova operação da ideia de distância.

A distância entre (2) e (6) é a mesma que a distância entre (6) e
(2). Ela não depende da direção.

Entretanto, os movimentos

\[
6\ominus2
\]

e

\[
2\ominus6
\]

não devem ser iguais.

O primeiro representa o movimento que parte de (2) e chega a (6).
Esse movimento aponta para a frente.

O segundo representa o movimento que parte de (6) e chega a (2).
Esse movimento aponta no sentido contrário.

Portanto, a operação que estamos a construir deve guardar não apenas o
comprimento do movimento, mas também a sua direção.

\subsection{Uma equação impossível nos naturais}

Consideremos agora a equação

\[
5+x=3.
\]

Nos números naturais, não existe um movimento para a frente que leve
(5) até (3).

Assim, a expressão

\[
3\ominus5
\]

não possui resultado em $(\mathbb{N}_0)$.

Diante dessa dificuldade, temos pelo menos duas possibilidades:

\begin{itemize}
\item declarar que a equação não possui solução em
$(\mathbb{N}_0)$;
\item ampliar o nosso conjunto numérico, introduzindo movimentos no
sentido contrário.
\end{itemize}

A primeira resposta não está errada. De facto, a equação

\[
5+x=3
\]

não possui solução em $(\mathbb{N}_0)$.

Mas podemos perguntar:

\begin{center}
\emph{Seria possível criar números capazes de representar movimentos
no sentido contrário?}
\end{center}

\subsection{Movimentos no sentido contrário}

O símbolo

\[
-2
\]

pode ser interpretado como um movimento de duas unidades no sentido
contrário ao movimento representado por (2).

Assim,

\[
5+(-2)=3.
\]

O número (-2) é chamado de \emph{oposto} de (2), pois os dois
movimentos anulam-se:

\[
2+(-2)=0.
\]

De maneira geral, para cada número natural (a), introduzimos um número
(-a) que satisfaz

\[
a+(-a)=0.
\]

O número (-a) é chamado de \emph{inverso aditivo} ou \emph{oposto} de
(a).

Ao juntar os números naturais, o zero e os seus opostos, obtemos o
conjunto dos números inteiros:

\[
\mathbb{Z}
=
{\ldots,-3,-2,-1,0,1,2,3,\ldots}.
\]

Na reta dos números inteiros, os números positivos representam
movimentos num sentido e os números negativos representam movimentos
no sentido contrário.

Agora a equação

\[
5+x=3
\]

possui solução em $\mathbb{Z}$:

\[
x=-2.
\]

\subsection{O reconhecimento da subtração}

A operação temporariamente representada por $(\ominus)$ é aquilo que
normalmente chamamos de \emph{subtração}.

Escrevemos

\[
a-b
\]

em vez de

\[
a\ominus b.
\]

A subtração pode ser definida usando a soma e o número oposto:

\[
a-b=a+(-b).
\]

Por exemplo,

\[
\begin{aligned}
6-2
&=6+(-2)\
&=4,
\end{aligned}
\]

enquanto

\[
\begin{aligned}
2-6
&=2+(-6)\
&=-4.
\end{aligned}
\]

A subtração não é uma operação comutativa, pois

\[
6-2\neq2-6.
\]

Isso responde a uma das perguntas deixadas anteriormente: existem
operações nas quais a ordem dos elementos altera o resultado.

\subsection{Os números negativos foram inventados ou descobertos?}

Essa pergunta não possui uma única resposta definitiva.

Os símbolos, as palavras e as regras usadas para representar os números
negativos foram inventados pelos seres humanos.

Entretanto, uma vez que desejamos resolver todas as equações da forma

\[
a+x=b,
\]

a necessidade de representar movimentos contrários surge naturalmente.
As relações que esses novos números precisam satisfazer não podem ser
escolhidas arbitrariamente.

Podemos, portanto, dizer que a linguagem foi inventada, mas a estrutura
e as suas consequências foram descobertas.

\subsection{Outras equações impossíveis}

A ampliação de $(\mathbb{N}_0)$ para $\mathbb{Z}$ não será a última.

Considere a equação

\[
2x=1.
\]

Ela não possui solução em $\mathbb{Z}$. Para resolvê-la, precisamos
introduzir números como

\[
\frac{1}{2}.
\]

Esses números conduzem ao conjunto dos números racionais:

\[
\mathbb{Q}.
\]

Considere agora

\[
x^2=2.
\]

Essa equação não possui solução em $\mathbb{Q}$. A sua solução positiva
é

\[
x=\sqrt{2},
\]

que é um número irracional. Os números racionais e irracionais formam o
conjunto dos números reais:

\[
\mathbb{R}.
\]

Por fim, a equação

\[
x^2=-1
\]

não possui solução em $\mathbb{R}$. A introdução de um número (i)
satisfazendo

\[
i^2=-1
\]

conduz ao conjunto dos números complexos:

\[
\mathbb{C}.
\]

Obtemos, assim, uma sequência de ampliações:

\[
\mathbb{N}_0
\subseteq
\mathbb{Z}
\subseteq
\mathbb{Q}
\subseteq
\mathbb{R}
\subseteq
\mathbb{C}.
\]

Cada ampliação permite resolver problemas que não possuíam solução no
conjunto anterior.

\section{Uma quantidade que pode variar}

\subsection{De incógnita a variável}

Até agora, usámos letras para representar valores desconhecidos.

Na equação

\[
2+x=6,
\]

a letra (x) representa um valor determinado que ainda não conhecemos.

Podemos, porém, usar uma letra de outra maneira.

Consideremos a regra

\[
y=x+2.
\]

Agora não estamos à procura de um único valor para (x). Permitimos que
(x) assuma diferentes valores e observamos o valor de (y)
correspondente.

Se (x=1), então

\[
y=1+2=3.
\]

Se (x=2), então

\[
y=2+2=4.
\]

Se (x=3), então

\[
y=3+2=5.
\]

Obtemos as correspondências

\[
1\mapsto3,\qquad
2\mapsto4,\qquad
3\mapsto5.
\]

A letra (x) pode variar, enquanto (y) é determinado pelo valor
escolhido para (x).

Por isso, chamamos (x) de \emph{variável independente} e (y) de
\emph{variável dependente}.

\subsection{O surgimento da função}

Uma \emph{função} é uma regra que associa a cada entrada permitida
exatamente uma saída.

Podemos chamar a regra anterior de (f) e escrever

\[
f(x)=x+2.
\]

Se desejamos permitir como entradas todos os números inteiros e obter
saídas inteiras, escrevemos

\[
f\colon\mathbb{Z}\longrightarrow\mathbb{Z},
\qquad
f(x)=x+2.
\]

O primeiro conjunto, $\mathbb{Z}$, é o conjunto das entradas
permitidas e é chamado de \emph{domínio}.

O segundo conjunto, também $\mathbb{Z}$, é o conjunto no qual as
saídas são procuradas e é chamado de $\emph{contradomínio}$.

Temos, por exemplo,

\[
f(1)=3,\qquad
f(2)=4,\qquad
f(3)=5.
\]

A escrita

\[
y=f(x)
\]

significa que (y) é a saída determinada pela função quando a entrada
é (x).

\subsection{Uma função que já utilizámos}

Quando introduzimos a indução matemática, usámos o símbolo (S(n)) para
representar o sucessor de (n).

Agora podemos reconhecer que (S) é uma função:

\[
S\colon\mathbb{N}_0\longrightarrow\mathbb{N}_0,
\]

definida por

\[
S(n)=n+1.
\]

Assim,

\[
S(0)=1,\qquad
S(1)=2,\qquad
S(2)=3.
\]

Já estávamos a usar uma função antes de termos dado um nome geral a esse
tipo de objeto.

\subsection{Diferença entre incógnita e variável}

Na equação

\[
2+x=6,
\]

a letra (x) é uma incógnita. Procuramos um valor específico que
satisfaça a equação.

Na função

\[
f(x)=x+2,
\]

a letra (x) é uma variável. Ela pode percorrer os valores permitidos
pelo domínio.

A diferença não está na letra usada, mas no papel que a letra
desempenha.

A mesma letra pode ser incógnita num problema e variável noutro.

Podemos responder às perguntas da seguinte maneira:

\begin{itemize}
\item Uma incógnita pode ser tratada como uma variável cujos valores
possíveis estão a ser testados, mas ela representa um valor fixo que
desejamos determinar.
\item Nem toda variável é uma incógnita, pois uma variável pode
percorrer livremente vários valores sem que desejemos descobrir um
único valor para ela.
\end{itemize}

\section{Movimentos que crescem regularmente}

\subsection{Uma regra de crescimento}

Consideremos a função

\[
f(x)=2x.
\]

Vamos observar os valores produzidos quando (x) percorre os números

\[
1,2,3,4.
\]

Temos

\[
\begin{aligned}
f(1)&=2\cdot1=2,\
f(2)&=2\cdot2=4,\
f(3)&=2\cdot3=6,\
f(4)&=2\cdot4=8.
\end{aligned}
\]

Obtemos a sequência

\[
2,4,6,8.
\]

Cada termo é obtido acrescentando duas unidades ao termo anterior:

\[
4-2=2,\qquad
6-4=2,\qquad
8-6=2.
\]

Dizemos que essa sequência possui uma \emph{diferença constante} igual
a (2).

Podemos representar o termo que ocupa a posição (n) por

\[
a_n=2n.
\]

Assim,

\[
a_1=2,\qquad
a_2=4,\qquad
a_3=6,\qquad
a_4=8.
\]

Uma sequência pode, portanto, ser vista como uma função cujo domínio é
formado pelos números naturais usados para indicar as posições.

Neste caso,

\[
a\colon\mathbb{N}\longrightarrow\mathbb{N},
\qquad
a(n)=2n.
\]

\subsection{Somando os movimentos}

Podemos agora somar os quatro primeiros valores da sequência:

\[
S_4=2+4+6+8.
\]

Usando a distributividade, obtemos

\[
\begin{aligned}
S_4
&=2\cdot1+2\cdot2+2\cdot3+2\cdot4\
&=2(1+2+3+4)\
&=2\cdot10\
&=20.
\end{aligned}
\]

A mesma ideia pode ser usada para os primeiros (n) números pares:

\[
S_n=2+4+6+\cdots+2n.
\]

Colocando o fator (2) em evidência,

\[
S_n=2(1+2+3+\cdots+n).
\]


Ou ainda usar a mesma idea para somar números ímpares, considere agora:

\[
1,
\]

\[
1+3,
\]

\[
1+3+5,
\]

\[
1+3+5+7.
\]

Calculando:

\[
1=1,
\]

\[
1+3=4,
\]

\[
1+3+5=9,
\]

\[
1+3+5+7=16.
\]

Obtemos:

\[
1,4,9,16,\ldots
\]

Esses números podem ser escritos como:

\[
1.1,\quad 2.2,\quad 3.3,\quad 4.4.
\]

Podemos introduzir uma nova notação:

\[
n^{\wedge 2}=n.n.
\]

Assim:

\[
n^{\wedge 2}=1+3+5+\cdots+(2n-1).
\]

Observe que a operação $n^{\wedge 2}$ é o mesmo que a de quadrado $n^2$. Fica a cargo do leitor fazer as demonstrações necessárias 
assim como fizemos para o $\mov$ e o $\rep$ para então chamarmos de quadrado. Assumimos que você o fez e a partir de agora Vamos
começar a usar a notação $n^2$ invés de $n^{\wedge 2}$.

O quadrado aparece como o resultado acumulado de movimentos ímpares.

\subsubsection{Por que os números ímpares formam quadrados?}

Começamos com um ponto:

\[
1.
\]

Para transformar um quadrado de lado 1 em um quadrado de lado 2, acrescentamos três pontos.

Para transformar um quadrado de lado 2 em um quadrado de lado 3, acrescentamos cinco pontos.

Para transformar um quadrado de lado 3 em um quadrado de lado 4, acrescentamos sete pontos.

Portanto:

\[
(n+1)^2=n^2+(2n+1).
\]

Perguntas:

\begin{itemize}
    \item Por que o incremento entre quadrados consecutivos é sempre ímpar?
    \item Que formas surgem ao somar números pares?
    \item Que formas surgem ao somar números triangulares?
\end{itemize}

Surge agora um novo problema:

\begin{center}
\emph{Como calcular a soma dos primeiros (n) números naturais sem
somá-los um por um?}
\end{center}

\subsection{A soma dos primeiros números naturais}

Definamos

\[
T_n=1+2+3+\cdots+(n-1)+n.
\]

Podemos escrever a mesma soma na ordem contrária:

\[
T_n=n+(n-1)+(n-2)+\cdots+2+1.
\]

Somando as duas expressões termo a termo, obtemos

\[
\begin{aligned}
2T_n
&=(1+n)+(2+n-1)+(3+n-2)+\cdots+(n+1).
\end{aligned}
\]

Cada par possui o mesmo valor:

\[
n+1.
\]

Como existem (n) pares, temos

\[
2T_n=n(n+1).
\]

Portanto,

\[
T_n=\frac{n(n+1)}{2}.
\]

A divisão por (2) sempre produz um número natural, pois entre dois
números consecutivos, (n) e (n+1), um deles é necessariamente par.

Voltando à soma dos números pares,

\[
\begin{aligned}
S_n
&=2T_n\
&=2\left(\frac{n(n+1)}{2}\right)\
&=n(n+1).
\end{aligned}
\]

Assim,

\[
2+4+6+\cdots+2n=n(n+1).
\]

Expandindo o produto,

\[
S_n=n^2+n.
\]

Começámos com uma sequência cujos termos crescem por movimentos
constantes:

\[
2,4,6,8,\ldots
\]

Ao acumular esses termos, surgiu uma expressão contendo $(n^2)$:

\[
n^2+n.
\]

Essa é a nossa primeira aparição natural de uma expressão polinomial.

\subsection{Do crescimento linear ao crescimento quadrático}

Os termos da sequência original são dados por

\[
a_n=2n.
\]

Eles crescem linearmente: de um termo para o seguinte, acrescentamos
sempre (2).

Entretanto, a soma acumulada desses termos é

\[
S_n=n^2+n.
\]

Ela não cresce acrescentando sempre a mesma quantidade. As quantidades
acrescentadas são justamente

\[
2,4,6,8,\ldots
\]

Temos

\[
S_{n+1}-S_n=2(n+1).
\]

Portanto, uma sequência de movimentos que cresce linearmente produz,
quando acumulada, uma quantidade que cresce quadraticamente.

Isso sugere novas perguntas:

\begin{itemize}
\item O que acontece quando somamos uma sequência com diferença
constante diferente de (2)?
\item Que tipo de expressão surge ao somarmos os quadrados
$1^2,2^2,3^2,\ldots$?
\item Se uma soma de termos lineares produz uma expressão
quadrática, o que produzirá uma soma de termos quadráticos?
\end{itemize}
\section{Quando a variação também começa a variar}

Consideremos novamente a função

\[
f(x)=2x.
\]

Quando \(x\) percorre os números

\[
1,2,3,4,
\]

obtemos

\[
\begin{aligned}
f(1)&=2\cdot1=2,\\
f(2)&=2\cdot2=4,\\
f(3)&=2\cdot3=6,\\
f(4)&=2\cdot4=8.
\end{aligned}
\]

Assim, a função produz a sequência

\[
2,4,6,8.
\]

A diferença entre termos consecutivos é sempre a mesma:

\[
4-2=2,
\]

\[
6-4=2,
\]

\[
8-6=2.
\]

Portanto, essa sequência possui diferença constante.

Consideremos agora as somas acumuladas dos seus termos:

\[
S_1=2,
\]

\[
S_2=2+4=6,
\]

\[
S_3=2+4+6=12,
\]

\[
S_4=2+4+6+8=20.
\]

Obtemos uma nova sequência:

\[
2,6,12,20.
\]

Vamos calcular as diferenças entre os seus termos:

\[
6-2=4,
\]

\[
12-6=6,
\]

\[
20-12=8.
\]

As diferenças são agora

\[
4,6,8.
\]

Elas já não são constantes. Portanto, a sequência

\[
2,6,12,20,\ldots
\]

não cresce linearmente.

Entretanto, as próprias diferenças formam uma nova sequência:

\[
4,6,8,\ldots
\]

E, ao calcularmos novamente as diferenças, encontramos

\[
6-4=2,
\]

\[
8-6=2.
\]

A primeira diferença não é constante, mas a diferença das diferenças é
constante.

Podemos organizar esse processo da seguinte forma:

\[
\begin{array}{ccccccccc}
2 && 6 && 12 && 20 && \cdots\\
& 4 && 6 && 8 &&\\
&& 2 && 2 &&
\end{array}
\]

A sequência original \(2,4,6,8,\ldots\) possui diferença constante logo
no primeiro nível.

A sequência acumulada \(2,6,12,20,\ldots\) só apresenta uma diferença
constante no segundo nível.

Isso sugere uma nova maneira de distinguir diferentes tipos de
crescimento:

\begin{center}
\emph{Quantas vezes precisamos calcular diferenças até obtermos uma
quantidade constante?}
\end{center}

A regra que produz a sequência acumulada é, portanto,

\[
S_n=n^2+n.
\]

Nessa expressão aparecem dois tipos diferentes de quantidade:

\[
n^2
\]

e

\[
n.
\]

O termo \(n\) já havia aparecido nas regras de crescimento linear.
Entretanto, o termo \(n^2\) representa uma nova ordem de crescimento.

Ao acumularmos os valores de uma regra linear, surgiu uma expressão
quadrática.

Isso conduz a novas perguntas:

\begin{itemize}
    \item Por que a soma de termos lineares produziu uma expressão
    quadrática?
    \item O que significam separadamente os movimentos \(n^2\) e \(n\)?
    \item Podemos somar movimentos de ordens diferentes?
    \item O que acontece quando somamos termos quadráticos?
    \item Surgirá então uma expressão contendo \(n^3\)?
\end{itemize}
\section{Podemos somar movimentos de ordens diferentes?}

Ao somarmos os primeiros termos da sequência

\[
2,4,6,8,\ldots,
\]

obtivemos os valores acumulados

\[
2,6,12,20,\ldots
\]

O termo que ocupa a posição \(n\) nessa nova sequência é dado por

\[
S_n=n^2+n.
\]

Pela primeira vez, encontramos numa mesma regra dois tipos diferentes
de movimento:

\[
n
\]

e

\[
n^2.
\]

O termo \(n\) representa um crescimento linear, enquanto \(n^2\)
representa um crescimento quadrático.

Isso sugere uma nova pergunta:

\begin{center}
\emph{Podemos construir regras somando movimentos de ordens diferentes?}
\end{center}

Por exemplo, podemos considerar a regra

\[
P(n)=n+n^2+n^3.
\]

Até agora, partíamos de uma sequência e procurávamos uma regra capaz de
produzi-la. Faremos agora o caminho inverso: começaremos com uma regra e
observaremos a sequência que ela produz.

Calculando os seus primeiros valores, obtemos

\[
\begin{aligned}
P(1)&=1+1^2+1^3=3,\\
P(2)&=2+2^2+2^3=14,\\
P(3)&=3+3^2+3^3=39,\\
P(4)&=4+4^2+4^3=84,\\
P(5)&=5+5^2+5^3=155,\\
P(6)&=6+6^2+6^3=258.
\end{aligned}
\]

Assim, a regra produz a sequência

\[
3,14,39,84,155,258,\ldots
\]

Vamos investigar como esses valores variam.

As primeiras diferenças são

\[
14-3=11,
\]

\[
39-14=25,
\]

\[
84-39=45,
\]

\[
155-84=71,
\]

\[
258-155=103.
\]

Obtemos, portanto,

\[
11,25,45,71,103,\ldots
\]

Essas diferenças não são constantes. Calculamos então as diferenças
novamente:

\[
25-11=14,
\]

\[
45-25=20,
\]

\[
71-45=26,
\]

\[
103-71=32.
\]

Obtemos

\[
14,20,26,32,\ldots
\]

Essas diferenças também não são constantes. Calculando uma terceira
vez:

\[
20-14=6,
\]

\[
26-20=6,
\]

\[
32-26=6.
\]

Finalmente, obtemos

\[
6,6,6,\ldots
\]

Podemos organizar o processo numa tabela de diferenças:

\[
\begin{array}{ccccccccccc}
3 && 14 && 39 && 84 && 155 && 258\\
& 11 && 25 && 45 && 71 && 103\\
&& 14 && 20 && 26 && 32\\
&&& 6 && 6 && 6
\end{array}
\]

Foi necessário calcular diferenças três vezes até encontrarmos uma
sequência constante.

Isso parece estar relacionado com a maior potência presente na regra:

\[
P(n)=n+n^2+n^3.
\]

A maior potência é \(n^3\), e a constância apareceu na terceira
diferença.

Para uma regra linear, como

\[
P(n)=2n,
\]

a primeira diferença é constante.

Para uma regra quadrática, como

\[
P(n)=n^2+n,
\]

a segunda diferença é constante.

Para a regra

\[
P(n)=n+n^2+n^3,
\]

a terceira diferença é constante.

Esses exemplos sugerem uma possível relação:

\begin{center}
\emph{O número de vezes que precisamos calcular diferenças até obter
uma constante parece coincidir com o maior expoente presente na regra.}
\end{center}

Entretanto, alguns exemplos não constituem uma demonstração.

Precisamos investigar:

\begin{itemize}
    \item Essa relação vale para toda soma de potências?
    \item O que acontece se multiplicarmos cada potência por um número?
    \item O que acontece se alguma potência estiver ausente?
    \item O termo de maior expoente determina sozinho a ordem das
    diferenças?
    \item Por que calcular uma diferença parece reduzir a maior potência?
\end{itemize}
\subsubsection{O que acontece se multiplicarmos cada potência por um número?}

Até agora, considerámos somas de movimentos como

\[
n+n^2+n^3.
\]

Entretanto, quando introduzimos a soma e a multiplicação, vimos que um
mesmo movimento pode aparecer mais de uma vez.

Por exemplo,

\[
3n
\]

pode ser interpretado como três ocorrências do movimento \(n\), enquanto

\[
4n^2
\]

representa quatro ocorrências do movimento \(n^2\).

Parece, portanto, natural permitir que cada ordem de movimento apareça
numa quantidade diferente. Podemos considerar, por exemplo,

\[
3+2n+4n^2+n^3.
\]

De maneira mais geral, podemos escrever

\[
a_1n+a_2n^2+a_3n^3,
\]

em que \(a_1,a_2,a_3\) são, inicialmente, números inteiros.

Também podemos acrescentar um termo que não depende de \(n\):

\[
a_0+a_1n+a_2n^2+a_3n^3.
\]

O termo \(a_0\) representa um movimento constante. O termo \(a_1n\)
representa um movimento de primeira ordem, \(a_2n^2\) representa um
movimento de segunda ordem e \(a_3n^3\) representa um movimento de
terceira ordem.

Os números

\[
a_0,a_1,a_2,a_3
\]

indicam quanto de cada ordem de movimento participa da regra. Esses
números serão chamados de \emph{coeficientes}.

Uma expressão construída por uma soma finita de potências inteiras não
negativas de uma variável, multiplicadas por coeficientes, será chamada
de \emph{polinómio}.

De maneira geral, escrevemos

\[
P(x)
=
a_0+a_1x+a_2x^2+\cdots+a_mx^m.
\]

O maior expoente que aparece com coeficiente diferente de zero será
chamado de \emph{grau do polinómio}.

\subsubsection{Quando o movimento total se anula?}

Uma nova pergunta pode ser feita sobre essas regras:

\begin{center}
\emph{Para quais movimentos de entrada o movimento total produzido pelo
polinómio é nulo?}
\end{center}

Consideremos o polinómio

\[
P(x)=x^2-5x+6.
\]

Procuramos valores de \(x\) para os quais

\[
P(x)=0.
\]

Experimentando alguns valores, encontramos

\[
P(2)=2^2-5\cdot2+6=0
\]

e

\[
P(3)=3^2-5\cdot3+6=0.
\]

Portanto, \(x=2\) e \(x=3\) anulam o movimento produzido pelo
polinómio.

Chamaremos de \emph{raiz} ou \emph{movimento nulo} de um polinómio todo
valor da variável que torna o valor do polinómio igual a zero.

Entretanto, encontrar as raízes apenas por experimentação pode tornar-se
difícil. Precisamos investigar se a própria estrutura do polinómio
guarda informações sobre os seus movimentos nulos.

\subsubsection{Reconstruindo um polinómio a partir dos seus movimentos
nulos}

Sabemos que \(2\) e \(3\) anulam o polinómio

\[
x^2-5x+6.
\]

Vamos procurar uma maneira de construir uma expressão que se anule
nesses dois valores.

A expressão

\[
x-2
\]

anula-se quando \(x=2\), enquanto

\[
x-3
\]

anula-se quando \(x=3\).

Se multiplicarmos essas duas expressões, obtemos

\[
(x-2)(x-3).
\]

Quando \(x=2\), o primeiro fator é zero e, portanto, todo o produto é
zero.

Quando \(x=3\), o segundo fator é zero e novamente todo o produto é
zero.

Expandindo o produto, encontramos

\[
\begin{aligned}
(x-2)(x-3)
&=x^2-3x-2x+6\\
&=x^2-5x+6.
\end{aligned}
\]

Assim,

\[
x^2-5x+6=(x-2)(x-3).
\]

O polinómio pode ser construído pelo produto das expressões associadas
aos seus movimentos nulos.

Vamos investigar um caso geral. Suponhamos que \(r\) e \(s\) sejam dois
movimentos nulos. Então consideramos

\[
(x-r)(x-s).
\]

Expandindo,

\[
\begin{aligned}
(x-r)(x-s)
&=x^2-sx-rx+rs\\
&=x^2-(r+s)x+rs.
\end{aligned}
\]

Portanto, se

\[
x^2+bx+c=(x-r)(x-s),
\]

devemos ter

\[
r+s=-b
\]

e

\[
rs=c.
\]

No exemplo

\[
x^2-5x+6,
\]

temos

\[
r+s=5
\]

e

\[
rs=6.
\]

Os números \(2\) e \(3\) satisfazem essas duas condições:

\[
2+3=5
\]

e

\[
2\cdot3=6.
\]

Isso conduz às seguintes observações:

\begin{itemize}
    \item o termo sem \(x\) é o produto dos dois movimentos nulos;
    \item o oposto do coeficiente de \(x\) é a soma dos dois movimentos
    nulos;
    \item o polinómio pode ser reconstruído pelo produto dos fatores
    associados aos seus movimentos nulos.
\end{itemize}

\subsubsection{Um movimento nulo igual a zero}

Consideremos agora

\[
x^2-5x.
\]

Colocando \(x\) em evidência, obtemos

\[
x^2-5x=x(x-5).
\]

Para que o produto seja zero, pelo menos um dos fatores precisa ser
zero. Portanto,

\[
x=0
\]

ou

\[
x-5=0,
\]

o que fornece

\[
x=5.
\]

Assim, os movimentos nulos são

\[
0
\qquad\text{e}\qquad
5.
\]

O termo constante do polinómio é zero e, de facto,

\[
0\cdot5=0.
\]

Além disso,

\[
0+5=5,
\]

que é o oposto do coeficiente \(-5\) da parcela linear.

Esses exemplos sugerem que os movimentos nulos estão diretamente
relacionados com os fatores do polinómio.

\subsubsection{Quantos movimentos nulos pode haver?}

Nos exemplos anteriores, os polinómios possuíam grau \(2\) e encontrámos
dois movimentos nulos.

Isso sugere uma nova conjectura:

\begin{center}
\emph{O grau do polinómio parece limitar a quantidade de movimentos
nulos que ele pode possuir.}
\end{center}

Entretanto, precisamos ser cuidadosos.

O polinómio

\[
x^2+1
\]

não possui movimentos nulos reais, pois

\[
x^2+1>0
\]

para todo número real \(x\).

Mais adiante, ao ampliarmos novamente o conjunto numérico, encontraremos
movimentos nulos complexos para esse polinómio.

Por enquanto, podemos apenas conjecturar que um polinómio de grau \(m\)
não pode possuir mais do que \(m\) movimentos nulos diferentes.

Também precisaremos considerar o caso em que um mesmo movimento nulo
aparece repetidamente.

\section{Movimentos que coincidem}

Nos exemplos anteriores, encontrámos dois movimentos nulos diferentes.
Vamos agora investigar o que acontece quando dois movimentos nulos
coincidem.

Suponhamos que o único movimento nulo seja \(x=a\), mas que o fator
associado a ele apareça duas vezes.

Nesse caso, temos

\[
(x-a)(x-a)=(x-a)^2.
\]

Expandindo,

\[
\begin{aligned}
(x-a)^2
&=x^2-ax-ax+a^2\\
&=x^2-2ax+a^2.
\end{aligned}
\]

A raiz \(a\) participa duas vezes da construção do polinómio.

Podemos dizer que \(a\) é uma raiz de multiplicidade \(2\).

Se o mesmo fator aparecer três vezes, teremos

\[
(x-a)^3=(x-a)(x-a)(x-a).
\]

Expandindo,

\[
\begin{aligned}
(x-a)^3
&=(x^2-2ax+a^2)(x-a)\\
&=x^3-ax^2-2ax^2+2a^2x+a^2x-a^3\\
&=x^3-3ax^2+3a^2x-a^3.
\end{aligned}
\]

Agora a raiz \(a\) aparece três vezes e possui multiplicidade \(3\).

Para quatro ocorrências, temos

\[
(x-a)^4=(x-a)(x-a)(x-a)(x-a).
\]

Podemos escrever

\[
(x-a)^4=(x-a)^2(x-a)^2.
\]

Assim,

\[
\begin{aligned}
(x-a)^4
&=(x^2-2ax+a^2)(x^2-2ax+a^2)\\
&=x^4-4ax^3+6a^2x^2-4a^3x+a^4.
\end{aligned}
\]

De maneira geral,

\[
\underbrace{(x-a)(x-a)\cdots(x-a)}
_{n\text{ fatores}}
=
(x-a)^n.
\]

O expoente \(n\) indica quantas vezes o mesmo movimento nulo participa
da construção do polinómio.

\subsection{Os coeficientes das potências repetidas}

Observemos os coeficientes das expansões anteriores, ignorando
momentaneamente os sinais e as potências:

\[
\begin{array}{ccccccccc}
&&&&1\\[4pt]
&&&1&&1\\[4pt]
&&1&&2&&1\\[4pt]
&1&&3&&3&&1\\[4pt]
1&&4&&6&&4&&1
\end{array}
\]

Os coeficientes parecem possuir uma organização simétrica.

Na expansão de \((x-a)^4\), por exemplo, encontramos

\[
1,4,6,4,1.
\]

Isso conduz a uma nova pergunta:

\begin{center}
\emph{Como podemos prever os coeficientes de \((x-a)^n\) sem realizar
todas as multiplicações?}
\end{center}

\subsection{De quantas maneiras uma parcela pode surgir?}

Para compreender de onde vêm os coeficientes, consideremos
temporariamente

\[
(x+a)^4
\]

em vez de

\[
(x-a)^4.
\]

Dessa maneira, podemos investigar primeiro a quantidade de parcelas sem
nos preocuparmos com a alternância dos sinais.

Temos

\[
(x+a)^4=(x+a)(x+a)(x+a)(x+a).
\]

Para formar uma parcela da expansão, precisamos escolher um elemento de
cada fator.

Em cada fator existem duas possibilidades:

\[
x
\qquad\text{ou}\qquad
a.
\]

Se escolhermos \(x\) nos quatro fatores, obtemos

\[
x\cdot x\cdot x\cdot x=x^4.
\]

Essa parcela só pode surgir de uma maneira, pois precisamos escolher
\(x\) em todos os fatores.

Se escolhermos \(a\) uma vez e \(x\) nos outros três fatores, obtemos
uma parcela da forma

\[
ax^3.
\]

Entretanto, podemos escolher o elemento \(a\) em qualquer um dos quatro
fatores:

\[
\begin{array}{cccc}
a & x & x & x,\\
x & a & x & x,\\
x & x & a & x,\\
x & x & x & a.
\end{array}
\]

Cada uma dessas escolhas produz a mesma parcela \(ax^3\).

Como existem quatro maneiras, encontramos

\[
4ax^3.
\]

Para obter uma parcela da forma

\[
a^2x^2,
\]

precisamos escolher \(a\) em dois fatores e \(x\) nos outros dois.

As possibilidades são

\[
\begin{array}{cccc}
a & a & x & x,\\
a & x & a & x,\\
a & x & x & a,\\
x & a & a & x,\\
x & a & x & a,\\
x & x & a & a.
\end{array}
\]

Existem seis maneiras de produzir a parcela \(a^2x^2\). Por isso,
encontramos

\[
6a^2x^2.
\]

Para obter uma parcela da forma

\[
a^3x,
\]

precisamos escolher \(x\) uma vez e \(a\) nos outros três fatores.

A posição do \(x\) pode ser escolhida de quatro maneiras. Portanto,
obtemos

\[
4a^3x.
\]

Finalmente, escolhendo \(a\) nos quatro fatores, obtemos

\[
a^4.
\]

Essa parcela só pode surgir de uma maneira.

Reunindo todas as parcelas, encontramos

\[
(x+a)^4
=
x^4+4ax^3+6a^2x^2+4a^3x+a^4.
\]

Os coeficientes contam quantas escolhas diferentes produzem a mesma
parcela.

\subsection{Um novo problema de contagem}

No caso anterior, precisámos contar quantas maneiras existem de escolher
duas posições para o elemento \(a\) entre quatro posições disponíveis.

Surge, então, um problema mais geral:

\begin{center}
\emph{De quantas maneiras podemos escolher \(p\) posições entre \(n\)
posições disponíveis?}
\end{center}

Chamaremos essa quantidade de

\[
C(n,p).
\]

Assim,

\[
C(4,2)=6,
\]

pois existem seis maneiras de escolher duas posições entre quatro.

Também temos

\[
C(4,1)=4
\]

e

\[
C(4,3)=4.
\]

Existem ainda dois casos imediatos.

Há somente uma maneira de não escolher posição alguma:

\[
C(n,0)=1.
\]

Também há somente uma maneira de escolher todas as posições:

\[
C(n,n)=1.
\]

\subsection{Construindo uma contagem a partir de outra}

Vamos investigar como calcular \(C(n,p)\) usando contagens menores.

Considere novamente a escolha de \(p\) posições entre \(n\) posições.
Observemos apenas a última posição.

Há duas possibilidades.

Na primeira, não escolhemos a última posição. Nesse caso, ainda
precisamos escolher \(p\) posições entre as \(n-1\) posições anteriores.

A quantidade de possibilidades é

\[
C(n-1,p).
\]

Na segunda possibilidade, escolhemos a última posição. Como uma posição
já foi escolhida, precisamos escolher apenas mais \(p-1\) posições entre
as \(n-1\) anteriores.

A quantidade de possibilidades é

\[
C(n-1,p-1).
\]

Como os dois casos são diferentes e juntos contêm todas as
possibilidades, temos

\[
C(n,p)=C(n-1,p)+C(n-1,p-1).
\]

Por exemplo,

\[
C(4,2)=C(3,2)+C(3,1).
\]

Como

\[
C(3,2)=3
\]

e

\[
C(3,1)=3,
\]

obtemos

\[
C(4,2)=3+3=6.
\]

\subsection{O triângulo das escolhas}

Usando

\[
C(n,0)=C(n,n)=1
\]

e a relação

\[
C(n,p)=C(n-1,p)+C(n-1,p-1),
\]

podemos construir sucessivamente as linhas

\[
\begin{array}{ccccccccccc}
 & & & & &1& & & & & \\[4pt]
 & & & &1& &1& & & & \\[4pt]
 & & &1& &2& &1& & & \\[4pt]
 & &1& &3& &3& &1& & \\[4pt]
 &1& &4& &6& &4& &1& \\[4pt]
1& &5& &10& &10& &5& &1
\end{array}
\]

Cada número interno é obtido pela soma dos dois números que aparecem
imediatamente acima dele.

Essa organização é tradicionalmente chamada de \emph{triângulo de
Pascal}.

As suas linhas fornecem exatamente os coeficientes das expansões de

\[
(x+a)^n.
\]

\subsection{Uma notação para contar escolhas}

A quantidade \(C(n,p)\) é tradicionalmente representada por

\[
\binom{n}{p}.
\]

Lemos essa expressão como ``\(n\) escolhe \(p\)''.

Assim,

\[
C(n,p)=\binom{n}{p}.
\]

A relação encontrada anteriormente torna-se

\[
\binom{n}{p}
=
\binom{n-1}{p}
+
\binom{n-1}{p-1}.
\]

Também temos

\[
\binom{n}{0}=1
\]

e

\[
\binom{n}{n}=1.
\]

Essa notação não representa uma divisão. Ela representa o número de
maneiras de escolher \(p\) posições entre \(n\) posições disponíveis.

\subsection{Uma fórmula direta para as escolhas}

O triângulo permite calcular uma linha usando a linha anterior.
Entretanto, podemos perguntar:

\begin{center}
\emph{É possível calcular diretamente o número de escolhas sem construir
todas as linhas anteriores?}
\end{center}

Suponhamos que desejamos escolher \(p\) posições entre \(n\) posições.

Se a ordem das escolhas fosse importante, teríamos \(n\) possibilidades
para a primeira posição.

Depois de escolhermos a primeira, restariam \(n-1\) possibilidades para
a segunda.

Depois, restariam \(n-2\) possibilidades para a terceira, e assim
sucessivamente.

Para escolher \(p\) posições em ordem, teríamos

\[
n(n-1)(n-2)\cdots(n-p+1)
\]

possibilidades.

Entretanto, quando escolhemos posições para formar uma parcela, a ordem
das escolhas não importa.

Escolher primeiro a posição \(1\) e depois a posição \(3\) produz o
mesmo conjunto de posições que escolher primeiro a posição \(3\) e
depois a posição \(1\).

Precisamos, portanto, descobrir quantas vezes cada conjunto de \(p\)
posições foi contado.

As mesmas \(p\) posições podem ser ordenadas de

\[
p(p-1)(p-2)\cdots2\cdot1
\]

maneiras.

Consequentemente,

\[
\binom{n}{p}
=
\frac{
n(n-1)(n-2)\cdots(n-p+1)
}{
p(p-1)(p-2)\cdots2\cdot1
}.
\]

Por exemplo,

\[
\binom{4}{2}
=
\frac{4\cdot3}{2\cdot1}
=
6.
\]

\subsection{Uma nova operação de compressão}

Produtos decrescentes como

\[
n(n-1)(n-2)\cdots2\cdot1
\]

aparecem repetidamente.

Para simplificar a escrita, definimos

\[
n!
=
n(n-1)(n-2)\cdots2\cdot1.
\]

Essa operação é chamada de \emph{fatorial}.

Por exemplo,

\[
4!=4\cdot3\cdot2\cdot1=24
\]

e

\[
3!=3\cdot2\cdot1=6.
\]

Também definimos

\[
0!=1.
\]

Podemos interpretar \(0!\) como um produto sem fatores. O valor \(1\)
permite conservar as propriedades das fórmulas que envolvem fatoriais.

Usando a nova notação, temos

\[
n(n-1)\cdots(n-p+1)
=
\frac{n!}{(n-p)!}.
\]

Portanto,

\[
\binom{n}{p}
=
\frac{n!}{p!(n-p)!}.
\]

A fórmula com fatoriais não cria o número de escolhas. Ela apenas
comprime a contagem que construímos anteriormente.

\section{O binómio de Newton}

Consideremos novamente

\[
(x+a)^n.
\]

Para obter uma parcela da forma

\[
x^{n-p}a^p,
\]

precisamos escolher \(a\) em exatamente \(p\) dos \(n\) fatores.

Nos \(n-p\) fatores restantes, escolhemos \(x\).

O número de maneiras de escolher as \(p\) posições ocupadas por \(a\) é

\[
\binom{n}{p}.
\]

Portanto, a parcela

\[
x^{n-p}a^p
\]

aparece

\[
\binom{n}{p}
\]

vezes na expansão.

Reunindo todas as possibilidades, obtemos

\[
\begin{aligned}
(x+a)^n
={}&
\binom{n}{0}x^n
+
\binom{n}{1}x^{n-1}a
+
\binom{n}{2}x^{n-2}a^2\\
&+\cdots+
\binom{n}{n-1}xa^{n-1}
+
\binom{n}{n}a^n.
\end{aligned}
\]

Usando o símbolo de soma para comprimir a expressão, podemos escrever

\[
(x+a)^n
=
\sum_{p=0}^{n}
\binom{n}{p}
x^{n-p}a^p.
\]

Essa identidade é chamada de \emph{binómio de Newton}.

Para obter a expansão de \((x-a)^n\), basta substituir \(a\) por
\(-a\):

\[
(x-a)^n
=
\sum_{p=0}^{n}
\binom{n}{p}
x^{n-p}(-a)^p.
\]

Como

\[
(-a)^p=(-1)^pa^p,
\]

temos

\[
(x-a)^n
=
\sum_{p=0}^{n}
(-1)^p
\binom{n}{p}
x^{n-p}a^p.
\]

Os sinais alternam porque cada escolha de \(a\) corresponde, na
realidade, à escolha de um fator \(-a\).

\subsection{Por que a expansão termina?}

A fórmula dos coeficientes pode ser escrita como

\[
\binom{n}{p}
=
\frac{
n(n-1)(n-2)\cdots(n-p+1)
}{p!}.
\]

Quando \(n\) é um número natural e \(p>n\), o produto do numerador
contém o fator

\[
n-n=0.
\]

Por exemplo,

\[
\binom{3}{4}
=
\frac{3\cdot2\cdot1\cdot0}{4!}
=
0.
\]

Todos os coeficientes posteriores também são zero.

Por essa razão, a expansão de

\[
(1+x)^n
\]

termina depois de um número finito de parcelas.

Ela é um polinómio.

\subsection{O expoente precisa ser um número natural?}

As expressões

\[
(1+x)^{-1}
\]

e

\[
(1+x)^{1/2}
\]

também possuem significado.

No primeiro caso,

\[
(1+x)^{-1}=\frac{1}{1+x}.
\]

No segundo,

\[
(1+x)^{1/2}=\sqrt{1+x}.
\]

Entretanto, os expoentes \(-1\) e \(1/2\) não representam uma quantidade
natural de fatores iguais.

Podemos perguntar:

\begin{center}
\emph{Seria possível conservar o padrão dos coeficientes do binómio
quando o expoente não é um número natural?}
\end{center}

Para um número qualquer \(\alpha\), podemos experimentar a definição

\[
\binom{\alpha}{0}=1
\]

e, para \(p\geq1\),

\[
\binom{\alpha}{p}
=
\frac{
\alpha(\alpha-1)(\alpha-2)\cdots(\alpha-p+1)
}{p!}.
\]

Quando \(\alpha=n\) é um número natural, recuperamos os coeficientes já
conhecidos.

Mas, quando

\[
\alpha=-1
\]

ou

\[
\alpha=\frac12,
\]

nenhum dos fatores precisa tornar-se zero.

A expansão pode, portanto, não terminar.

Isso sugere a expressão

\[
(1+x)^\alpha
\stackrel{?}{=}
\sum_{p=0}^{\infty}
\binom{\alpha}{p}x^p.
\]

O símbolo de interrogação indica que ainda não sabemos se essa igualdade
é verdadeira.

Antes de aceitá-la, precisamos responder a novas perguntas:

\begin{itemize}
    \item O que significa somar uma quantidade infinita de parcelas?
    \item As somas parciais aproximam-se de algum valor determinado?
    \item Para quais valores de \(x\) a aproximação funciona?
    \item O erro pode tornar-se menor do que qualquer quantidade
    previamente escolhida?
    \item Podemos somar, multiplicar, derivar e acumular essas expansões
    como fazemos com polinómios finitos?
\end{itemize}

Essas perguntas conduzem ao estudo das \emph{séries infinitas}.


\end{document}